\section{The spike-corrected SVM}
\label{sec:scsvm}

Corollary~\ref{cor:1} shows that the SVM does not hold the consistency property
under the spiked model. The interpretation in Section~\ref{sec:example} gives a
remedy: a spiked direction is a source of noise which cannot be learned when the
sample size is fixed, so that we should remove it if it carries no information about
the classification.

\begin{definition}
\label{def:1}
For a direction $h_{is}$, we define the signal-to-spike ratio by
\[
\eta_{is}=\frac{(h_{is}^T\mu)^2}{\lambda_{is}}=\frac{\Delta b_{is}^2}{\lambda_{is}}
\longrightarrow\frac{\delta\beta_{is}^2}{c_{is}}.
\]
\end{definition}

Under~(A3) and~(S), it holds that $\eta_{is}\to\delta\beta_{is}^2/c_{is}$.

\begin{proposition}
\label{prop:2}
Consider the subproblem obtained by projecting the data onto a one-dimensional
direction $h$. When $n_i$ is fixed, the limiting misclassification rate obtained by
using this direction is not smaller than $\Phi(-\sqrt{\eta}/2)$. Hence, when
$\eta_{is}\to\infty$, the direction should be kept, and when $\eta_{is}=O(1)$, the
direction should be removed.
\end{proposition}

\begin{proof}[Outline of proof]
Projecting onto $h$, the population $\pi_i$ becomes a one-dimensional distribution
with mean $h^T\mu_i$ and variance $\lambda$. The mean difference is $|h^T\mu|$ and
the standard deviation is $\sqrt{\lambda}$, so that the error rate of the Bayes
rule with known parameters is $\Phi\{-|h^T\mu|/(2\sqrt{\lambda})\}=\Phi(-\sqrt{\eta}/2)$
in the normal case. When $\eta=O(1)$, this is bounded below by a positive constant,
and the estimation error makes it larger. On the other hand, when $\eta\to\infty$,
it tends to zero. As for the loss caused by the removal, we note that
$\eta_{is}=O(1)$ implies $(h_{is}^T\mu)^2=O(\lambda_{is})=O(\theta)$, so that the
loss is negligible when the residual signal diverges faster than $\theta$.
\end{proof}

\subsection{The procedure}

\begin{definition}[SC-SVM]
\label{def:2}
We propose the spike-corrected SVM as follows.
\begin{enumerate}
\item Split the observations from each $\pi_i$ into $D_i^{(1)}$ for estimation and
$D_i^{(2)}$ for training. The split is essential in the theory because it makes
the estimated projection independent of the training data.
\item Apply the NR methodology of \citet{yata2012} to $D_i^{(1)}$ and obtain the
number of the spikes $\hat m_i$, the eigenvalues $\tilde\lambda_{is}$ and the
eigenvectors $\hat h_{is}$.
\item Compute $\hat\eta_{is}=(\hat h_{is}^T\hat\mu)^2/\tilde\lambda_{is}$ with
$\hat\mu=\bar x_1^{(1)}-\bar x_2^{(1)}$, and collect the directions with
$\hat\eta_{is}\le\tau_d$ into the set $\hat S$, where $\tau_d\to\infty$ is a
threshold such as $\tau_d=\log d$.
\item Let $\hat U$ be an orthonormal basis of $\hat S$ and put
$\hat P=I_d-\hat U\hat U^T$. Project the training data and the new observation
onto $\hat P$.
\item Apply the BC-SVM with the modified bias term $\hat\kappa_*$ given in
Proposition~\ref{prop:1} to the projected data. When some spiked directions are
kept, construct a plug-in rule on the low-dimensional part and add it to the
discriminant function.
\end{enumerate}
\end{definition}

We assume the following conditions.

\begin{enumerate}[label=(C\arabic*),leftmargin=*]
\item It holds that $n_i\to\infty$ and $n_i/d\to 0$.
\item For $s\le m_i$, it holds that $\tilde\lambda_{is}/\lambda_{is}\to 1$ and
$\hat h_{is}^Th_{is}\to 1$ in probability.
\item Let $P_*$ be the projection onto the orthogonal complement of the removed
spikes and $\Delta_*=\|P_*\mu\|^2$. It holds that
\[
\Delta_*/\{\tr(\Sigma_{1*}^2)+\tr(\Sigma_{2*}^2)\}^{1/2}\to\infty.
\]
\end{enumerate}

The condition~(C2) holds under the conditions given by \citet{yata2012}, such as
$n_i\tr(\Sigma_{i*}^2)/\lambda_{im_i}^2\to 0$. We assume it explicitly in this
paper.

\begin{theorem}
\label{thm:5}
Assume \textup{(A1)} to \textup{(A4)}, \textup{(S)} and \textup{(C1)} to
\textup{(C3)}. Then the SC-SVM holds the consistency property, that is,
\[
e_{SC}(1)+e_{SC}(2)\longrightarrow 0
\quad\text{in probability as }d\to\infty.
\]
On the other hand, from Corollary~\ref{cor:1}, the SVM and the BC-SVM do not hold
the consistency property under the same conditions.
\end{theorem}

\begin{proof}
(i) We first consider the ideal procedure in which $\hat P$ is replaced by the true
projection $P_*$. The eigenvalues of $\Sigma_i^P=P_*\Sigma_i P_*$ consist of
$\{\lambda_{is}:s\notin\text{removed}\}$ and the kept spikes. For the part from
which the noise spikes are removed, we have
\[
\frac{\tr\{(\Sigma_i^P)^2\}}{\{\tr(\Sigma_i^P)\}^2}\to 0
\]
from~(S). That is, the projected data meet the condition~\eqref{eq:1.2} of the
non-spiked model.

(ii) Under the non-spiked model, the geometric representation is recovered, so that
the consistency of the BC-SVM given by \citet{nakayama2017} is applicable under
(C3). We use the modified bias term $\kappa_*$ given in Proposition~\ref{prop:1}.

(iii) We replace $P_*$ by $\hat P$. From~(C2) we have $\|\hat P-P_*\|=o_P(1)$.
Since $\hat P$ is independent of the training data and $x_0$ from the sample
splitting, we may argue conditionally. Noting that the spiked coordinates of $x$
are of order $O_P(\sqrt{\theta})$, we have
\[
\|(\hat P-P_*)x\|=o_P(\sqrt{\theta}),
\]
so that the contribution to the inner products is $o_P(\theta)$. Hence it is
absorbed into the remainder terms of Lemmas~\ref{lem:1} to~\ref{lem:3} and the
conclusions of~(i) and~(ii) are preserved.

(iv) From~(C2) and $n_i\to\infty$, the quantity $\hat\eta_{is}$ is relatively
consistent for $\eta_{is}$. Choosing $\tau_d\to\infty$ such that
$\tau_d=o(\eta_{is})$ for the directions to be kept, the selection is correct with
probability tending to one.
\end{proof}
